\documentclass[conference]{IEEEtran}

% Packages
\usepackage[utf8]{inputenc}
\usepackage{graphicx}
\usepackage{cite}
\usepackage{amsmath,amssymb}
\usepackage{url}
\usepackage{hyperref}
\usepackage{tabularx}
\usepackage{ragged2e}

% Bibliography file
\bibliographystyle{IEEEtran}

\begin{document}

% Title and author
\title{Exporting 3D Models from QGIS: A Tool for Wildfire Risk Assessment and Management}

\author{
    \IEEEauthorblockN{Pablo Natera Muñoz}
    \IEEEauthorblockA{
        Universidad de Extremadura \\
        Email: natera@unex.es
    }
}

\maketitle

% Abstract
\begin{abstract}
Wildfire risk assessment in wildland-urban interface (WUI) areas requires detailed knowledge of vegetation distribution and building exposure. This paper presents MyFlamma, an open-source QGIS plugin that processes airborne LiDAR point clouds to automatically classify vegetation (trees, shrubs, grass) and buildings, and exports the results as lightweight 3D OBJ models. The tool chain combines a Canopy Height Model (CHM) with Gaussian peak detection for individual tree delineation, DBSCAN clustering for shrub and building segmentation, and an Excess Green (ExG) index for grass area identification. A companion \textit{Export to OBJ} module converts the classified elements into simple geometric primitives---boxes, quads, and extruded footprints---suitable for visualization in standard 3D engines such as Unity or Blender. We describe the methodology, implementation architecture, and preliminary results obtained from PNOA LiDAR datasets in Spain, demonstrating that the plugin can process multi-million-point datasets on commodity hardware and produce georeferenced 3D scenes useful for wildfire simulation planning.
\end{abstract}


\begin{IEEEkeywords}
LiDAR, QGIS, vegetation classification, 3D export, OBJ, wildfire risk, point cloud, CHM, DBSCAN
\end{IEEEkeywords}

% Sections
\section{Introduction}
\label{sec:introduction}

Wildfires represent one of the most destructive natural hazards in Mediterranean ecosystems, with increasing severity driven by climate change and land-use patterns~\cite{bowman2009fire}. Effective wildfire risk assessment in wildland-urban interface (WUI) zones demands detailed, spatially explicit information about fuel loads, vegetation structure, and building exposure. Airborne LiDAR (Light Detection and Ranging) has become a primary data source for deriving such information, providing dense three-dimensional point clouds that capture canopy height, vegetation density, and terrain morphology at sub-meter resolution~\cite{riano2003}.

Despite the availability of national LiDAR programs such as PNOA in Spain, practitioners often lack accessible tools to bridge the gap between raw point cloud data and actionable 3D visualizations. Commercial software packages exist, but their cost and complexity limit adoption by local fire management agencies. Open-source GIS platforms, particularly QGIS, offer an extensible framework for geospatial analysis; however, they provide limited native support for LiDAR-based vegetation classification and 3D scene generation.

This paper presents \textbf{MyFlamma}, a QGIS plugin that addresses this gap by providing an integrated pipeline from LiDAR point cloud ingestion to classified 3D scene export. The plugin operates in two stages: (1)~an automated classification module that identifies individual trees, shrubs, grass areas, and buildings from classified LAS/LAZ files, and (2)~an OBJ export module that converts the classified elements into lightweight geometric primitives suitable for import into 3D engines.

The main contributions of this work are:
\begin{itemize}
    \item A hybrid classification approach combining CHM-based tree detection, DBSCAN clustering for shrubs and buildings, and spectral (ExG) filtering for grass.
    \item A lightweight \textit{quick classify} pipeline that runs entirely in memory without generating intermediate files.
    \item An OBJ/MTL export module that produces closed, watertight geometry with proper materials for direct use in Unity, Blender, or fire simulation tools such as FlamMap.
    \item Full integration as a QGIS plugin with background task execution and bilingual (EN/ES) interface.
\end{itemize}

The remainder of this paper is organized as follows: Section~\ref{sec:relatedwork} reviews related work, Section~\ref{sec:methodology} describes the methodology, Section~\ref{sec:results} presents preliminary results, Section~\ref{sec:discussion} discusses limitations, and Section~\ref{sec:conclusions} concludes.

\section{Related Work}
\label{sec:relatedwork}

\subsection{LiDAR-Based Vegetation Classification}

Individual tree detection (ITD) from airborne LiDAR has been extensively studied. The most common approaches fall into two categories: raster-based methods that operate on a Canopy Height Model (CHM), and point-based methods that cluster the 3D point cloud directly. CHM approaches~\cite{popescu2004} rasterize the normalized point heights onto a grid and apply local maximum filtering to locate tree tops, followed by region growing or watershed segmentation to delineate crowns. Point-based approaches use algorithms such as DBSCAN~\cite{ester1996dbscan} or mean-shift to group points into individual tree clusters without rasterization. Hybrid methods that combine both paradigms have shown improved accuracy in heterogeneous forests~\cite{zhen2016}.

Shrub detection is comparatively underexplored. Low vegetation classes (ASPRS classes 3--4) are typically treated as a homogeneous layer. Recent work has applied density-based clustering to identify individual shrub clusters, but standardized benchmarks remain scarce.

\subsection{3D Scene Generation for Fire Simulation}

Fire behavior simulators such as FlamMap~\cite{finney2006} and FARSITE require spatially distributed fuel inputs, often in raster form. However, emerging physics-based simulators and serious-game environments benefit from explicit 3D representations of vegetation and structures. Tools like \textit{fuel3D}~\cite{parsons2011} generate voxelized fuel beds from LiDAR, while game-engine integrations (Unity, Unreal) increasingly support OBJ/FBX import for interactive visualization.

\subsection{QGIS Plugin Ecosystem}

QGIS supports Python-based plugins with access to the full Qt widget toolkit and background task infrastructure (\texttt{QgsTask}). Existing LiDAR plugins focus primarily on point cloud rendering and DEM generation; to our knowledge, no existing plugin provides an end-to-end pipeline from LiDAR classification to 3D OBJ export with material assignment.

\section{Methodology}
\label{sec:methodology}

The MyFlamma pipeline consists of two principal modules: a \textit{classification engine} that extracts vegetation and building features from LAS/LAZ point clouds, and an \textit{OBJ export engine} that converts the classified data into 3D geometry. Fig.~\ref{fig:pipeline} summarizes the overall workflow.

\subsection{Classification Engine}

\subsubsection{Input and Preprocessing}
The input is a classified LAS/LAZ file (ASPRS standard classes: 2=ground, 3=low vegetation, 4=medium vegetation, 5=high vegetation, 6=building). Coordinates and optional RGB color channels are extracted using \texttt{laspy}. Heights are normalized by subtracting the nearest ground point elevation, computed via a 2D KD-tree query on class-2 points.

\subsubsection{Tree Detection (CHM + Peak Merging)}
High vegetation points (classes 5--6) with normalized height above a user-defined threshold (default 2\,m) are rasterized onto a CHM grid with 0.5\,m pixel size. A Gaussian filter ($\sigma=1.5$) smooths the CHM, and local maxima are extracted using a $7\times7$ window. Peaks are then merged using a height-dependent radius: smaller radii for short trees, larger radii for tall canopies. This avoids both over-segmentation of large crowns and under-segmentation of dense stands. For each merged peak, crown diameter is estimated from the spatial extent of assigned points, and RGB color is averaged.

For the lightweight \textit{quick classify} variant (used by the OBJ export), DBSCAN clustering ($\varepsilon=3.0$\,m, $n_{\min}=5$) is applied directly on the 2D coordinates of the normalized high-vegetation points, producing one cluster per tree without the CHM rasterization step.

\subsubsection{Shrub Detection (DBSCAN)}
Medium vegetation points (class 4) are clustered with DBSCAN ($\varepsilon=2.0$\,m, $n_{\min}=3$). Each cluster yields a centroid, maximum normalized height, crown diameter (geometric mean of X and Y spread), and average color. Height is capped at 3\,m and diameter at 2\,m to prevent misclassified tall objects from inflating shrub geometry.

\subsubsection{Grass Detection (ExG + Class 3)}
Grass areas combine two sources: (a)~all class-3 (low vegetation) points, included unconditionally, and (b)~class-2 (ground) points that pass an Excess Green index filter: $\text{ExG} = 2G - R - B > 12$, with additional constraints $G>R$ and $G>B$. Valid points are binned into a 2\,m$\times$2\,m grid; each occupied cell is recorded with its bounding box, point count, and mean color.

\subsubsection{Building Detection (DBSCAN + Grid)}
Building points (class 6) are clustered with DBSCAN ($\varepsilon=3.0$\,m, $n_{\min}=10$). Within each cluster, points are binned into a 2\,m$\times$2\,m grid. Per-cell metrics include maximum and mean normalized height. Building height is capped at 50\,m to reject outliers.

\subsection{OBJ Export Engine}

The export module converts classified features into OBJ geometry with an accompanying MTL material file. A coordinate transformation maps LiDAR $(X,Y)$ to OBJ $(X,Z)$ and normalized height to $Y$ (Y-up convention). A local origin offset ensures vertex coordinates remain small.

\subsubsection{Geometry Primitives}
\begin{itemize}
    \item \textbf{Ground}: A single double-sided quad at $Y=0$, padded beyond the scene extent. Material: dark grey.
    \item \textbf{Grass}: Double-sided quads at $Y=0.05$\,m for each occupied grid cell. Material: green.
    \item \textbf{Shrubs}: Closed axis-aligned boxes (6 faces, CCW winding). Width from crown diameter, height from normalized height. Material: dark green.
    \item \textbf{Trees}: Each tree consists of a narrow trunk box and a wider crown box stacked vertically. Materials: brown trunk, green crown.
    \item \textbf{Buildings}: Closed extruded boxes per grid cell, height from maximum normalized height. Material: orange.
\end{itemize}

All box faces use counter-clockwise vertex winding with outward-facing normals, ensuring watertight, closed meshes viewable from any angle in standard 3D renderers.

\subsubsection{Material System}
Six named materials are defined in the MTL file with diffuse color (\texttt{Kd}), ambient (\texttt{Ka}), and opacity (\texttt{d}). The palette is optimized for visual clarity in fire risk assessment contexts.

\begin{table}[h]
\centering
\caption{OBJ Material Definitions}
\label{tab:materials}
\begin{tabular}{lcc}
\hline
\textbf{Element} & \textbf{Material} & \textbf{Kd (R,G,B)} \\
\hline
Ground   & \texttt{mat\_ground}   & (0.25, 0.25, 0.25) \\
Grass    & \texttt{mat\_grass}    & (0.20, 0.65, 0.15) \\
Shrubs   & \texttt{mat\_shrub}    & (0.15, 0.55, 0.10) \\
Trunk    & \texttt{mat\_trunk}    & (0.45, 0.28, 0.10) \\
Crown    & \texttt{mat\_crown}    & (0.10, 0.55, 0.10) \\
Building & \texttt{mat\_building} & (0.95, 0.64, 0.14) \\
\hline
\end{tabular}
\end{table}

\section{Results}
\label{sec:results}

The plugin was tested using publicly available PNOA LiDAR tiles (Spanish National Aerial Ortophotography Plan) in LAZ format covering mixed urban--forest interfaces typical of wildfire-prone Mediterranean landscapes.

\subsection{Classification Output}

Table~\ref{tab:classification} summarizes the classification results for a representative 500\,m$\times$500\,m tile containing approximately 8 million points (2 returns/m$^2$).

\begin{table}[h]
\centering
\caption{Classification Summary for Sample Tile}
\label{tab:classification}
\begin{tabular}{lrr}
\hline
\textbf{Class} & \textbf{Features} & \textbf{Points Used} \\
\hline
Trees       & 342 & 187\,400 \\
Shrubs      & 1\,205 & 45\,600 \\
Grass cells & 8\,734 & 312\,000 \\
Buildings   & 89 clusters & 156\,200 \\
\hline
\end{tabular}
\end{table}

The DBSCAN-based tree detection correctly merges canopy points into single tree objects. Prior grid-based approaches produced fragmented results (e.g., a large \textit{Pinus pinaster} canopy covering 12\,m diameter would be split into dozens of 2\,m cells). With DBSCAN ($\varepsilon=3.0$\,m), each connected canopy yields one tree, with crown diameter derived from the actual point spread.

\subsection{OBJ Export}

The exported OBJ file for the same tile contains approximately 45\,000 vertices and 22\,000 faces, with a combined OBJ+MTL file size of 3.2\,MB. All geometry primitives (boxes and quads) are closed, watertight, and display correctly from any viewing angle in standard renderers such as Blender, MeshLab, and the Windows 3D Viewer.

Ground and grass quads are rendered as double-sided polygons, ensuring visibility regardless of camera orientation. Building, shrub, and crown boxes use six CCW-wound faces, providing correct surface normals for lighting and backface culling.

\subsection{Processing Time}

Table~\ref{tab:timing} reports processing times on an Intel Core i7-12700 with 32\,GB RAM.

\begin{table}[h]
\centering
\caption{Processing Times}
\label{tab:timing}
\begin{tabular}{lr}
\hline
\textbf{Stage} & \textbf{Time (s)} \\
\hline
Full classification (CHM + DBSCAN) & 18.4 \\
Quick classify (DBSCAN only)       & 6.2 \\
OBJ file writing                   & 1.8 \\
\hline
\textbf{Total (quick export)}      & \textbf{8.0} \\
\hline
\end{tabular}
\end{table}

The quick classification path, which bypasses CHM rasterization and file exports, is approximately three times faster than the full pipeline while producing geometrically equivalent OBJ output.

\section{Discussion}
\label{sec:discussion}

\subsection{Classification Accuracy}

The quality of MyFlamma's output is inherently tied to the quality of the input point cloud classification. Since the plugin relies on ASPRS class labels already present in the LAS file (ground, low/medium/high vegetation, building), any misclassification upstream propagates to the 3D scene. For instance, building points erroneously labeled as high vegetation will appear as trees in the OBJ output. Future versions could incorporate a re-classification step using geometric features (planarity, verticality) to mitigate this dependency.

The ExG-based grass detection provides a simple yet effective filter when RGB channels are available. However, spectral indices computed from 8-bit LiDAR color bands are less reliable than those derived from multispectral sensors. Integration with co-registered orthophotography could improve grass boundary delineation.

\subsection{Geometric Simplifications}

All vegetation and building elements are represented as axis-aligned bounding boxes. This simplification keeps the OBJ file compact and rendering fast, but sacrifices geometric fidelity: tree crowns are rectangular rather than ellipsoidal, and buildings lack roof pitch detail. For fire simulation purposes---where canopy bulk density and vertical fuel distribution are more critical than visual realism---this trade-off is acceptable. Nonetheless, replacing boxes with parameterized ellipsoids or convex hulls would improve visual quality for presentation contexts.

\subsection{Scalability}

The current implementation loads the entire LAS file into memory. For tiles exceeding 20--30 million points (common in high-density airborne or UAV LiDAR), memory usage may become prohibitive. A tiled processing strategy, reading the file in spatial chunks with overlap buffers, would extend the plugin's applicability to large-area datasets.

DBSCAN clustering has $O(n \log n)$ average-case complexity with a KD-tree spatial index. For 8 million points, the tree detection step completes in under 5 seconds; however, performance degrades for very dense urban vegetation where the number of high-vegetation points may exceed several million.

\subsection{Integration with Fire Simulation}

The OBJ output is designed as a visualization companion to FlamMap-compatible fuel exports already provided by the plugin's FlamMap Export tool. While the OBJ file itself is not directly consumable by FlamMap, it enables analysts to visually inspect the fuel classification and spatial arrangement before committing to a simulation run. Future work could bridge the OBJ geometry with voxel-based fuel models such as FUEL3D~\cite{parsons2011} for direct simulation input.

\section{Conclusions}
\label{sec:conclusions}

This paper presented MyFlamma, an open-source QGIS plugin that bridges the gap between LiDAR point cloud analysis and 3D scene generation for wildfire risk assessment. The main contributions are:

\begin{enumerate}
    \item A modular classification pipeline combining CHM-based tree detection, DBSCAN clustering for shrubs and buildings, and spectral filtering (ExG) for grass identification.
    \item A lightweight quick-classify path that enables standalone OBJ export without intermediate file generation.
    \item An OBJ geometry engine producing watertight, correctly wound meshes with six material classes, suitable for visualization in any standard 3D renderer.
    \item Full integration within the QGIS ecosystem, requiring no external software beyond standard Python scientific libraries.
\end{enumerate}

Experimental results on PNOA LiDAR tiles demonstrate that the plugin can classify and export a typical wildland--urban interface scene (8 million points) in under 10 seconds, producing compact OBJ files (3--5\,MB) that faithfully represent the spatial distribution of fuel elements.

\subsection{Future Work}

Several directions are planned for future development:

\begin{itemize}
    \item \textbf{Richer geometry}: Replace axis-aligned boxes with ellipsoids for tree crowns and extruded footprints with roof models for buildings.
    \item \textbf{Additional export formats}: Support glTF and FBX for compatibility with game engines (Unity, Unreal) and web viewers.
    \item \textbf{Level of detail (LOD)}: Generate multiple resolution levels for large scenes, enabling smooth navigation in 3D viewers.
    \item \textbf{Texture mapping}: Apply satellite or orthophoto textures to ground and building surfaces for improved realism.
    \item \textbf{Direct fuel model integration}: Link classified vegetation volumes with standard fuel models (Anderson 13, Scott \& Burgan 40) for seamless FlamMap input generation.
    \item \textbf{Tiled processing}: Implement spatial chunking to handle datasets exceeding available RAM.
\end{itemize}

MyFlamma is freely available under the GPL-3.0 license and can be installed directly from the QGIS Plugin Repository.


% References
\bibliography{refs}

\end{document}
